\chapter{Conclusion}\label{ch:conclusion}

Four months of work on elliptic-curve index calculus produced four kinds of
result.

\paragraph{Theory.}
\begin{itemize}
\item \emph{What is fixed for every curve.} The summation cover's monodromy
  is $(\ZZ/2)^{m-2}$ everywhere. Its fibres split completely where relation
  search operates. The mean yield at a given base size is a double count.
  So no choice of curve or base geometry can raise the \emph{mean}
  decomposition yield at matched base size. Coverage, the cost of the
  membership oracle and the difficulty of the solve remain free, and the
  engineering of Part~\ref{part:algorithms} works on those.
\item \emph{What is open.} For three presentations of the decomposition
  system the first fall degree is a constant of the presentation, and the
  last fall degree is not bounded by it:
  \begin{itemize}
  \item proved for the guarded chain;
  \item certified on named instances of binary Weil descent;
  \item measured for prime-field Kummer towers.
  \end{itemize}
  The subexponential claims for binary curves assume exactly the step
  these results show is not automatic. Whether the last fall degree grows
  without bound is unproved.
\end{itemize}

\paragraph{Engineering.} The decomposition oracle became two to five
orders of magnitude cheaper where root finding applies. The Gröbner
engines improved by one to two orders in counted work. Factor bases
closed under known-logarithm maps, a $k$-sum table, and set-up and linear
algebra engineering each cut a constant. None of this changed an exponent
on a prime field or a prime-degree binary field. The prime-field $k$-sum
oracle moved its exponent from $0.62$ to $0.52$, but it is generic, and
its limit is baby-step giant-step.

\paragraph{Measurement.} At a matched reference, priced across a band and
replayed, rho is cheaper on every prime-order curve the programme measured:
\begin{itemize}
\item $4.7$--$18\times$ on Koblitz curves of 37--47 bits;
\item $5.3$--$102\times$ on prime curves of 20--40 bits;
\item $10.6$--$38\times$ on random binary curves.
\end{itemize}
Nineteen leaderboard rows passed a validation gate that demands independent
replay, recomputation and a matched rho. None of them is below its
reference. Every online win disappeared under an equal-advice audit, and
the corrections forced by matching, pricing and replaying lengthened the
measured distance more often than any single improvement shortened it.

\paragraph{Where index calculus wins.} In weak extension-field families
($E_0/\FF_{q^3}$ through genus-3 covers, Joux--Vitse covers of
$E(\FF_{p^6})$, and perhaps $k=4$ in characteristic two) the measured cost
crosses or approaches rho, as theory predicts. The plane-quartic route
measured exponent $\approx1$ against rho's $1.5$, at $0.005\times$ negation
rho near $q=2^{21}$. These are known algorithms; the programme's
contribution there is fair measurement.

\paragraph{What would change the picture.} Three items, in order of what
they would settle:
\begin{enumerate}
\item \emph{A proved lower bound on the last fall degree that grows with
  $n$}, for any presentation of the binary decomposition system. It would
  answer the question left open by Huang, Kosters and Yeo and by
  Kousidis and Wiemers, and it is the only result that bears on the
  subexponential claims.
\item \emph{More certified instances} at $n=33$--$51$, with public
  certificates, to show whether the last fall degree reaches 6.
\item \emph{The frozen $k=4$ replication} at $m=14$--$17$, the cheapest
  experiment that could put a measured crossover on record. It would be
  in a weak family.
\end{enumerate}
Nothing here concerns a deployed curve. The open problems of
Chapter~\ref{sec:open} state each remaining question as the experiment or
proof that would settle it.
